% SPDX-License-Identifier: AGPL-3.0-or-later
% Copyright (C) 2026 Gaurav Jadhav <gauravmakarandjadhav@gmail.com>
% Canonical source: https://github.com/GgauravJ05/kokoro
%
% Compile with: pdflatex kokoro.tex (twice, for references)
% or paste into Overleaf. CI builds this automatically; see .github/workflows/paper.yml
\documentclass[11pt,a4paper]{article}

\usepackage[utf8]{inputenc}
\usepackage[T1]{fontenc}
\usepackage[margin=1in]{geometry}
\usepackage{amsmath,amssymb}
\usepackage{booktabs}
\usepackage{graphicx}
\usepackage{hyperref}
\usepackage{xcolor}
\usepackage{caption}

\hypersetup{colorlinks=true, linkcolor=blue!50!black, urlcolor=blue!50!black,
            citecolor=blue!50!black}

\title{\textbf{Kokoro}: Mood-Conditioned Cold-Start Retrieval for Anime,\\
       and a Falsified Trajectory Hypothesis}
\author{Gaurav Jadhav\\
        \small\texttt{gauravmakarandjadhav@gmail.com}\\
        \small\url{https://github.com/GgauravJ05/kokoro}}
\date{September 2026}

\begin{document}
\maketitle

\begin{abstract}
Catalog metadata does not express mood. The genre \emph{Drama} contains both
\emph{Violet Evergarden} and \emph{School Days}, and no tag separates them on
the axis a viewer actually cares about. We learn an affective representation of
anime from user review prose instead, using a two-tower contrastive retriever
trained with InfoNCE on (review segment $\rightarrow$ title metadata) pairs.
Because the item tower encodes metadata only, the model can place a title that
has never been interacted with. On a cold-start split of $9{,}864$ titles this
matters absolutely rather than marginally: every collaborative baseline we test
scores exactly $0.0000$ NDCG@10, because a model whose only representation of an
item is who interacted with it has no representation of an item nobody has
interacted with. Our trained content retriever reaches $0.0845$ NDCG@10 under
the standard cold-only protocol, $4.97\times$ random and $+41\%$ over an
off-the-shelf sentence encoder, served at $4.4$\,ms p95.

We additionally report two negative results that we consider part of the
contribution. First, routing the item tower through eight named bipolar mood
axes with an anchor-alignment loss demonstrably shapes the space --- it triples
axis spread and corrects two inverted axes at no retrieval cost --- but only
three of eight axes reach even weak agreement (AUC $0.63$--$0.69$) with held-out
tag probes, and none clears our stated $0.70$ bar. Second, and centrally, we
falsify our own motivating hypothesis: that a title is better modelled as a mood
\emph{curve} over narrative position than as a point. Curves can be built, and a
permutation control confirms a real position-dependent signal ($1.15\times$ the
movement of position-shuffled curves), but the curve's \emph{shape} carries no
information beyond its own mean --- shape-only prediction sits at chance
($0.485$ AUC) and adding shape features degrades a mean-only model
($0.513$ vs.\ $0.547$). We report the conditions under which this conclusion
should not be generalised.
\end{abstract}

\section{Introduction}

Recommender systems for narrative media are usually built on two signals:
collaborative filtering over interaction histories, and content filtering over
catalog metadata. Both fail on the query a viewer actually asks. ``Something
melancholy but hopeful, slow, that won't wreck me'' is a statement about
affect, and affect is not in the metadata.

Two observations motivate this work.

\textbf{Mood must be learned from language.} MyAnimeList and AniList expose
genres, themes, studios and scores. None distinguishes works that occupy
opposite emotional territory inside one genre label. But viewers write at length
about exactly this, in reviews, and that text is a supervision signal.

\textbf{Interaction-based models cannot serve new titles.} A seasonal catalog
adds work continuously, and a model that represents an item by its interaction
column has, by construction, nothing to say about an item with an empty column.
This is not a degradation; it is undefined.

We take the second observation as the evaluation frame. Our headline split is
cold-start, and the claim we defend is narrow and falsifiable: a content tower
trained contrastively on review prose can place titles no collaborative model
can represent at all.

We also set out to defend a third claim --- that mood is a trajectory rather
than a point --- and failed. Section~\ref{sec:trajectory} reports that failure
with its controls, because a negative result with a permutation control is more
informative than the weak positive we could have produced by tuning until the
numbers cooperated.

\section{Data}

Both live APIs this project was designed against became unavailable during the
work: AniList returned \texttt{403 The AniList API has been temporarily disabled
due to severe stability issues}, and Jikan returned \texttt{504} for any
endpoint requiring a live MyAnimeList fetch. Ingestion was therefore made
source-agnostic and now defaults to static public dumps, which has the secondary
benefit of exact reproducibility without an account or API key.

\begin{table}[h]
\centering
\caption{Corpus composition. Nothing is redistributed; \texttt{data/} holds only
ingestion scripts and a manifest recording each input's SHA-256, measured join
rate, licence and known biases.}
\begin{tabular}{lrl}
\toprule
Component & Size & Source (licence) \\
\midrule
Catalog titles        & $28{,}880$    & \texttt{lyfesan/myanimelist-top-anime} (MIT) \\
Descriptive tags      & $4{,}736$ titles & \texttt{LeData/media-metadata-anilist} (CC0) \\
Reviews               & $52{,}520$    & \texttt{Raiser1/anime-review-mal} (other) \\
User--item ratings    & $6{,}332{,}233$ & CooperUnion MAL dump (undeclared) \\
Rated titles          & $9{,}864$     & --- \\
Users                 & $69{,}599$    & --- \\
\bottomrule
\end{tabular}
\end{table}

\subsection{Constraints we consider load-bearing}

\textbf{Supervision covers 1.7\% of the catalog.} Reviews span $449$ titles
against $28{,}880$. The model is supervised on the head and must generalise
through content, which is precisely why cold-start is the honest headline rather
than a footnote.

\textbf{No synopsis exists.} The catalog source has no synopsis field. An early
version of our pipeline matched a column named \texttt{Synonyms} with a
substring heuristic and silently trained the item tower on alternate titles. We
substitute AniList descriptive tags (\emph{Tragedy}, \emph{Philosophy},
\emph{Found Family}), which are closer to mood vocabulary than genres, at
$47.2\%$ coverage of rated titles.

\textbf{No interaction timestamps.} The rating matrix carries none, so a
temporal split is not computable and we do not report one. Debut dates are
available from the catalog, which is what makes a genuine cold-start split
possible.

\textbf{Review text is lowercased} and drawn from MyAnimeList's top-500 with
$91.3\%$ positive sentiment and \emph{mixed} verdicts dropped upstream. These
bound what any affective axis can learn and are recorded in the corpus manifest.

\section{Method}

\subsection{Two-tower contrastive retrieval}

A query tower encodes free text; an item tower encodes a title's metadata
string. Both project a frozen MiniLM sentence encoder into a shared
$256$-dimensional space, and we train only the projections. This is deliberate:
the off-the-shelf baseline is the identical model with identity projections, so
any difference is attributable to the objective rather than to encoder capacity.

Training uses InfoNCE with in-batch negatives,
\[
\mathcal{L} = -\frac{1}{B}\sum_{i=1}^{B}
  \log \frac{\exp(q_i^\top v_i / \tau)}
            {\sum_{j=1}^{B}\exp(q_i^\top v_j / \tau)},
\]
symmetrised over both directions, with $\tau = 0.05$.

Three design decisions carry the result.

\textbf{The item tower never sees reviews.} A cold title has none. An item tower
trained on review text would have nothing to encode for exactly the titles the
evaluation is about.

\textbf{Titles are masked out of queries.} A review naming its own work lets the
model match on the title string rather than on mood --- a shortcut invisible in
training metrics and fatal at cold start.

\textbf{Each batch contains a title at most once.} With $194$ trainable titles,
uniform sampling repeats titles constantly, and InfoNCE treats every off-diagonal
entry as a negative; a repeat is therefore a gradient actively separating two
segments that describe the same work.

\subsection{Named mood axes}
\label{sec:axes}

Optionally the item representation passes through a bottleneck of eight bipolar
axes (\emph{comfort}, \emph{pace}, \emph{hope}, \emph{intimacy},
\emph{cognition}, \emph{levity}, \emph{catharsis}, \emph{intensity}), squashed
to $[-1,1]$ and expanded back to the retrieval width. An anchor-alignment hinge
pulls each axis's natural-language pole phrases toward that pole:
\[
\mathcal{L}_{\text{anchor}} =
  \frac{1}{|A|}\sum_{a \in A} \max\!\big(0,\; 0.8 - p_a \cdot z_{a, k_a}\big),
\]
where $z$ is the bottleneck output, $k_a$ the anchor's axis and
$p_a \in \{-1,+1\}$ its pole.

\section{Evaluation}

\subsection{Splits}

We report \texttt{user\_holdout} for in-catalog accuracy and \texttt{cold\_start}
as the headline. A temporal split is not computable on this data and is not
reported. Cold-start is evaluated under two protocols, because reporting only
one is how such results get overstated: \emph{full catalog} (rank new titles
against everything, realistic and dominated by warm distractors) and
\emph{cold-only} (rank new titles against each other, the standard protocol).

\subsection{A metric that is misleading on this split}

Popularity lift is standard in beyond-accuracy reporting and is
\emph{uninterpretable} on a cold-start split. Every correct answer there has
zero training popularity by construction, so the oracle lift is $0.00\times$ and
a model achieves a \emph{high} lift precisely by filling slots with warm items
that cannot be correct. We therefore introduce \textbf{candidate share@$k$}: the
fraction of the top-$k$ that is eligible to be correct at all. It states the
same fact without the misreading.

\section{Results}

\subsection{In-catalog retrieval}

\begin{table}[h]
\centering
\caption{\texttt{user\_holdout}, $k{=}10$, $3{,}000$-user subsample. The
from-scratch NumPy BPR implementation wins on accuracy \emph{and} surfaces
$2.9\times$ more of the catalog at lower popularity bias.}
\begin{tabular}{lrrrr}
\toprule
Model & Recall@10 & NDCG@10 & Coverage@10 & Pop.\ lift \\
\midrule
Random          & 0.0010 & 0.0026 & 0.991 & 0.99$\times$ \\
Popularity      & 0.1047 & 0.1530 & 0.007 & 23.2$\times$ \\
Item-$k$NN      & 0.1630 & 0.2434 & 0.057 & 18.8$\times$ \\
\textbf{BPR-MF} & \textbf{0.1762} & \textbf{0.2633} & \textbf{0.167} & \textbf{13.6}$\times$ \\
\bottomrule
\end{tabular}
\end{table}

\subsection{Cold start}

\begin{table}[h]
\centering
\caption{Cold start over $1{,}491$ titles with zero training interactions
($998$ carry held-out ratings). Both content rows use the identical backbone and
user-profile logic; the only difference is the trained projection.}
\begin{tabular}{lrrrrr}
\toprule
& \multicolumn{3}{c}{Cold-only protocol} & \multicolumn{2}{c}{Full catalog} \\
\cmidrule(lr){2-4}\cmidrule(lr){5-6}
Model & Recall@10 & NDCG@10 & vs.\ random & NDCG@10 & Cand.\ share \\
\midrule
Popularity        & 0.0016 & 0.0032 & 0.19$\times$ & \textbf{0.0000} & \textbf{0.000} \\
Item-$k$NN        & 0.0016 & 0.0032 & 0.19$\times$ & \textbf{0.0000} & \textbf{0.000} \\
BPR-MF            & 0.0062 & 0.0147 & 0.87$\times$ & \textbf{0.0000} & \textbf{0.000} \\
Random            & 0.0078 & 0.0170 & 1.00$\times$ & 0.0017 & 0.149 \\
Content (off-the-shelf) & 0.0274 & 0.0597 & 3.51$\times$ & 0.0112 & 0.064 \\
\textbf{Content (trained)} & \textbf{0.0487} & \textbf{0.0845} & \textbf{4.97}$\times$
  & \textbf{0.0165} & \textbf{0.109} \\
\bottomrule
\end{tabular}
\end{table}

Every collaborative model is at or below random. Popularity and item-$k$NN reach
$0.19\times$ because they assign identical scores to all cold items and return
one constant list to every user (coverage@10 $=0.001$, Gini $=0.999$).
Contrastive training improves NDCG@10 by $41\%$ over the off-the-shelf encoder
and raises candidate share by $69\%$.

\subsection{Mood axes: a mechanism that works, a claim that does not}

\begin{table}[h]
\centering
\caption{Anchor-weight ablation. Agreement is ROC-AUC against held-out tag
probes; the $21$ probe tags are withheld from the item text during training, so
the model never reads \emph{Tragedy} on a title whose \emph{comfort} it is
graded on.}
\begin{tabular}{lrrrr}
\toprule
Anchor weight & Val recall@1 & Axis spread $\sigma$ & Mean AUC & Inverted axes \\
\midrule
0.0 (ablation) & 0.175 & 0.085 & 0.465 & \textbf{2} \\
1.0            & 0.183 & 0.186 & 0.537 & 0 \\
5.0            & 0.176 & 0.218 & 0.561 & 0 \\
20.0           & 0.179 & 0.238 & \textbf{0.565} & 0 \\
\bottomrule
\end{tabular}
\end{table}

The mechanism works and is nearly free: axis spread roughly triples, both
inverted axes are corrected, and retrieval moves within noise. Per axis,
\emph{levity} reaches $0.686$, \emph{comfort} $0.663$ and \emph{hope} $0.630$;
\emph{cognition}, \emph{intensity} and \emph{intimacy} show no signal, and
\emph{pace} and \emph{catharsis} have no tag proxy anywhere in $348$ AniList
tags and are reported as untested rather than assumed.

No axis clears our stated $0.70$ bar. The supportable claim is therefore
\emph{three of eight axes are weakly but measurably aligned with their names},
not that the model is interpretable. We note the confound: a failing probe is
not proof of a failing axis, and our \emph{intimacy} probe arguably tests cast
structure rather than loneliness.

\subsection{Mood trajectories: a falsified hypothesis}
\label{sec:trajectory}

This project was motivated by the argument that a title is a curve over
narrative position. \emph{Steins;Gate} inverts at episode 12; a single mood
vector averages its halves into something describing neither. We tested that
argument and it did not survive.

\textbf{Curves can be built, at coarse resolution.} $8.0\%$ of review segments
carry a resolvable position, a median of $194$ per title, yielding $5$-arc curves
for $140$ titles. But $62\%$ of positions come from words like ``the opening''
and ``the finale'', only $8\%$ from an explicit episode number, and a
named-arc pattern fired $3$ times in $141{,}000$ segments because the source is
lowercased. The evidence supports beginning/middle/end, not per-episode detail.

\textbf{The movement is real but small.} Shuffling each title's mood vectors
across its own segments destroys any position--mood relationship while
preserving sample size, bucket sizes and the marginal distribution:

\begin{table}[h]
\centering
\begin{tabular}{lrrr}
\toprule
Statistic & Real curves & Position-shuffled & Ratio \\
\midrule
Within-title $\sigma$      & 0.0827 & 0.0720 & 1.15$\times$ \\
Mean arc-to-arc step       & 0.1075 & 0.0967 & 1.11$\times$ \\
Within $\div$ between      & 1.010  & 0.839  & 1.20$\times$ \\
\bottomrule
\end{tabular}
\end{table}

Real curves move $15\%$ more than shuffled ones, so a position-dependent signal
exists --- but roughly $87\%$ of the apparent movement survives shuffling.

\textbf{The shape carries nothing beyond the mean.} Predicting held-out AniList
tags by 5-fold cross-validated logistic regression:

\begin{table}[h]
\centering
\begin{tabular}{lr}
\toprule
Features & Mean AUC \\
\midrule
Mean only (what a point model has)        & \textbf{0.547} \\
Shape only (curve minus its own mean)     & 0.485 \emph{(chance)} \\
Mean $+$ shape                            & 0.513 \emph{(worse than mean alone)} \\
\bottomrule
\end{tabular}
\end{table}

Shape-only prediction is at chance and adding shape degrades the model; only
$2$ of $9$ tags improved. We therefore conclude that \textbf{on this corpus the
trajectory hypothesis is not supported}, and we did not train the neural
trajectory encoder --- fitting a higher-capacity model to a signal that fails
its permutation control would be fitting noise.

Three conditions bound this conclusion and we do not offer them as excuses:
$140$ titles against $40$ shape features (the full-vs-mean loss is textbook
overfitting); three-region position resolution, which cannot express ``inverts
at episode 12''; and tag prediction possibly being an unfair target, since tags
describe whole works. A fair test needs case-sensitive, episode-level discussion
text --- for example per-episode discussion threads --- which we could not
obtain.

\subsection{Serving}

\begin{table}[h]
\centering
\caption{Latency on a laptop CPU, $9{,}864$ titles, $256$-dim vectors. HNSW at
\texttt{ef\_search}$=64$ achieves recall@10 $=1.000$ against exact search.}
\begin{tabular}{lrr}
\toprule
Stage & p50 & p95 \\
\midrule
Query encoding (MiniLM)      & 2.78\,ms & 3.04\,ms \\
HNSW lookup                  & 0.06\,ms & 0.08\,ms \\
\textbf{End-to-end (HTTP)}   & \textbf{3.20\,ms} & \textbf{4.44\,ms} \\
\bottomrule
\end{tabular}
\end{table}

We note that the index barely earns its place at this scale: exact brute force
runs in $0.102$\,ms, so HNSW saves $\approx 0.04$\,ms while the encoder costs
$3$\,ms. The budget is dominated by one small encoder forward pass, which is the
concrete argument for this architecture over a generative model call per
request.

\section{Limitations}

Absolute cold-start numbers are low ($0.0845$ NDCG@10); this is a first result,
not a solved problem. Supervision covers $1.7\%$ of the catalog. The review
corpus is drawn from a top-500 subset with $91.3\%$ positive sentiment, so the
axes are better resolved on the positive pole. Tag probes are a proxy for human
mood judgements, which we did not collect; distinguishing ``the axis is
meaningless'' from ``the probe is bad'' requires them.

\section{Conclusion}

Learning mood from review prose lets a retriever place titles that
interaction-based models cannot represent at all --- the difference between
$0.0845$ and exactly $0.0000$ NDCG@10 --- at $4.4$\,ms p95. The interpretability
mechanism functions but supports only a weak claim, and the trajectory
hypothesis that motivated the project is not supported by this data. We report
all three outcomes because a system that is only evaluated where it succeeds has
not been evaluated.

\section*{Reproducibility}

All results are produced by the public repository at
\url{https://github.com/GgauravJ05/kokoro} under AGPL-3.0. The corpus is built
from public dumps by \texttt{kokoro corpus build} and pinned by a content hash
(\texttt{a11aac2a2c226689}); every artifact carries a provenance stamp. Tables
in this paper are regenerated by \texttt{scripts/} entry points named in the
repository README.

\end{document}
